\section{Implementing the first best}\label{sec:mkt:implementation}

\subsection{The policy}\label{sec:mkt:implementation:policy}

Section~\ref{sec:mkt:setup:equilibrium} ended with a count: three shadow prices of the planned economy
--- $p^P$ on the pollution stock, $p^W$ on the ledger and $\zeta$ on the intensity condition --- have
no private counterpart, and the instrument list contains exactly three independent rates. The
proposition below assigns one to each.

\begin{proposition}[Decentralization of the first best]\label{prop:mkt:implementation}
Let $\left\{p^P,p^W,\zeta,\Phi^I\right\}_{t\ge0}$ denote the shadow prices and the investment-goods
material intensity along the solution to Problem~\ref{prob:sp}, and suppose that solution is unique.
Set
\begin{align}
\tau^{\Xi} &= p^P,
&
\tau^{\mathcal D} &= \zeta,
&
\tau^{\circ} &= p^W,
\label{eq:mkt:implement-primitives}
\end{align}
so that, by~\eqref{eq:mkt:emission-split} and the deposit--refund structure of block~(B),
\begin{align}
\tau^U &= p^P,
&
\tau^T &= s^R = d^Wp^P,
&
\tau^R &= p^W-\zeta,
&
s^I &= \Phi^I\left(p^W-\zeta\right),
\label{eq:mkt:implement-derived}
\end{align}
with $s^W$ set by the tender and $\mathcal T$ balancing the budget. Then the competitive equilibrium of
Definition~\ref{def:mkt:equilibrium} implements the first-best allocation, and the equilibrium prices
are
\begin{align}
P^K &= F_K\left(1-\zeta\Omega^Y\right),
&
P^R &= \Psi-p^W,
&
P^T &= \alpha V,
&
s^W &= p^W,
\label{eq:mkt:implement-prices}
\end{align}
while the private shadow prices coincide with the social ones, $P^S=p^S$, $P^X=p^X$, $p^M$ as
in~\eqref{eq:sp:sum:M}, and $q=1-\Phi^I\left(1-\sigma\right)\left(p^W-p^M\right)$ as
in~\eqref{eq:sp:sum:I}.
\end{proposition}

\begin{proof}
We verify that the first-best path, together with the prices~\eqref{eq:mkt:implement-prices}, satisfies
every condition of Definition~\ref{def:mkt:equilibrium}; uniqueness of the planner's solution then
makes the equilibrium the first best rather than merely one equilibrium among several. Throughout,
substitute~\eqref{eq:mkt:implement-primitives}--\eqref{eq:mkt:implement-derived} for the policy rates,
so that $\tau^{\mathcal D}=\zeta$ replaces the gate fee wherever it appears. Two identities do most of
the work:
\begin{align}
\Psi &= P^R+\tau^R+\tau^{\mathcal D} \;=\; P^R+p^W,
&
V &= \Psi-p^W+d^Wp^P \;=\; P^R+s^R .
\label{eq:mkt:price-identities}
\end{align}
The first is the definition $\Psi\equiv F_R\left(1-\zeta\Omega^Y\right)+\zeta$
of~\eqref{eq:sp:shorthand} combined with the second of~\eqref{eq:mkt:final-foc}; the second follows
from it and $s^R=d^Wp^P$. In words: the market price of a tonne is the planner's gross material value
net of the waste charge the tonne will eventually attract, and the virgin displacement value is the
market price of material grossed up by the recovery subsidy.

\emph{Final goods.} The first of~\eqref{eq:mkt:final-foc} is the definition of $P^K$
in~\eqref{eq:mkt:implement-prices}. The second is the first of~\eqref{eq:mkt:price-identities},
which is~\eqref{eq:sp:shorthand} rearranged. Both hold identically.

\emph{Extraction and exploration.} Using $\tau^{\circ}=p^W$ and $P^S=p^S$, the first
of~\eqref{eq:mkt:mining-foc} reads $P^R=C^N_N\left(1+\zeta\Omega^N\right)+p^S+p^W\Omega^{N,S}$, which
by~\eqref{eq:mkt:price-identities} is $\Psi=C^N_N\left(1+\zeta\Omega^N\right)+p^S+p^W\left(1+\Omega^{N,S}\right)$:
the planner's~\eqref{eq:sp:sum:N}. The second of~\eqref{eq:mkt:mining-foc} is
$p^S+p^X=C^D_D\left(1+\zeta\Omega^D\right)+p^W\Omega^{D,S}$, the planner's~\eqref{eq:sp:sum:D},
with no tax on either margin. The corner branches coincide because the two sides do. Finally,
the mining costates~\eqref{eq:mkt:mining-costates} are literally the planner's~\eqref{eq:sp:sum:S}
and~\eqref{eq:sp:sum:X} once $\tau^{\mathcal D}=\zeta$, so $P^S=p^S$ and $P^X=p^X$ follow from the
common discount rate and the transversality conditions.

\emph{Recycling capital.} By the second of~\eqref{eq:mkt:price-identities}, $P^R+s^R=V$, so the
branching inequality of~\eqref{eq:mkt:recycler-K-corner}--\eqref{eq:mkt:recycler-K-interior} is
$a'\!\left(0^+\right)V\lessgtr F_K\left(1-\zeta\Omega^Y\right)$: this is exactly the Kuhn--Tucker
system~\eqref{eq:sp:sum:recycling-virgin-KR}, branch for branch.

\emph{Treated waste and the treatment share.} The recycler's demand
condition~\eqref{eq:mkt:recycler-T} gives $P^T=\alpha V$, which is the third
of~\eqref{eq:mkt:implement-prices}. The waste firm's margin is then
\begin{align*}
P^T+\tau^U-\tau^T &= \alpha V+p^P-d^Wp^P \;=\; \alpha V+p^P\left(1-d^W\right),
\end{align*}
and the marginal cost in~\eqref{eq:mkt:waste-foc} is
$c^{T\prime}\!\left(\varpi\right)\left(1+\zeta\Omega^W\right)$. Both sides therefore agree with the
planner's~\eqref{eq:sp:sum:recycling-virgin-varpi}, again branch for branch, including the two corners.

\emph{The collection subsidy.} Substituting $P^T=\alpha V$, $\tau^U=p^P$ and $\tau^T=d^Wp^P$
into the zero-profit condition~\eqref{eq:mkt:waste-zero-profit},
\begin{align*}
s^W &= c^W\left(1+\zeta\Omega^W\right)+p^P\left(1-\varpi\right)+d^Wp^P\varpi-\alpha\varpi V .
\end{align*}
The planner's waste condition~\eqref{eq:sp:sum:W}, with its pollution bracket expanded as
$p^P\left(1-\varpi\right)+d^Wp^P\varpi-d^Wp^P\alpha\varpi$, reads
$p^W\left(1-\alpha\varpi\right)=c^W\left(1+\zeta\Omega^W\right)+p^P\left(1-\varpi\right)+d^Wp^P\varpi
-\alpha\varpi\left(\Psi+d^Wp^P\right)$. Since $V=\Psi-p^W+d^Wp^P$, the two right-hand sides differ by
$\alpha\varpi p^W$ exactly, so $s^W=p^W\left(1-\alpha\varpi\right)+\alpha\varpi p^W=p^W$.

\emph{The household.} With $\tau^{\mathcal D}=\zeta$ the consumption rule
in~\eqref{eq:mkt:hh-foc} is the planner's~\eqref{eq:sp:sum:C}. With $\tau^{\circ}=p^W$ the private
law of motion for $p^M$ is the planner's~\eqref{eq:sp:sum:M}, so the two coincide under the common
transversality condition. Then, by the definition of $q$ in~\eqref{eq:mkt:hh-costate-def} and
$s^I=\Phi^I\left(p^W-\zeta\right)$,
\begin{align*}
q &= 1-\Phi^I\left(p^W-\zeta\right)+\Phi^Ip^M
\;=\; 1-\Phi^I\left(p^W-p^M-\zeta\right)
\;=\; 1-\Phi^I\left(1-\sigma\right)\left(p^W-p^M\right),
\end{align*}
the last step by $\zeta=\sigma\left(p^W-p^M\right)$ from~\eqref{eq:sp:zeta-explicit}. This is the
planner's~\eqref{eq:sp:sum:I}. The arbitrage condition in~\eqref{eq:mkt:hh-foc} with
$P^K=F_K\left(1-\zeta\Omega^Y\right)$ is then the planner's~\eqref{eq:sp:costate:K}, and the
Keynes--Ramsey rule~\eqref{eq:sp:sum:Cgrowth} follows as before.

\emph{Accounting and clearing.} The relations~\eqref{eq:mkt:accounting} are the ledger and the
intensity condition of Problem~\ref{prob:sp}, which the first-best path satisfies by construction; the
market-clearing conditions~\eqref{eq:mkt:clearing} are the definitions~\eqref{eq:sp:definitions} and
the goods constraint~\eqref{eq:sp:goods}. Every condition of Definition~\ref{def:mkt:equilibrium} is
therefore met.
\end{proof}

\begin{corollary}[The materials block is self-financing]\label{cor:mkt:budget}
Under the policy of Proposition~\ref{prop:mkt:implementation}, net government revenue at every date is
\begin{align}
\mathcal B &= p^P\,\Xi .
\label{eq:mkt:budget}
\end{align}
The entire materials apparatus --- both gate fees, the materials input tax, the investment rebate and
the collection subsidy --- balances exactly, and the only net revenue in the economy is the Pigouvian
emission tax.
\end{corollary}

\begin{proof}
Collecting the seven flows and using $W=\Omega\mathcal D+\delta M^K+\mathcal N$ from
\eqref{eq:sp:ledger-streams} and $\Omega\mathcal D=R-\Phi^IG$ from~\eqref{eq:sp:intensity},
\begin{align*}
\mathcal B &= \underbrace{\zeta\,\Omega\mathcal D+p^W\left(\delta M^K+\mathcal N\right)}_{\text{gate fees}}
+\underbrace{\left(p^W-\zeta\right)R-\Phi^I\left(p^W-\zeta\right)G}_{\text{deposit--refund}}
+\underbrace{p^P\Xi}_{\text{emissions}}
-\underbrace{p^WW}_{\text{collection}}
\\
&= \zeta\,\Omega\mathcal D+\left(p^W-\zeta\right)\Omega\mathcal D-p^W\Omega\mathcal D+p^P\Xi
\;=\;p^P\Xi ,
\end{align*}
the primary-waste terms cancelling against the corresponding part of $p^WW$.
\end{proof}

\subsection{Reading the policy}\label{sec:mkt:implementation:reading}

\smalltitle{Three prices, three instruments.} The assignment
in~\eqref{eq:mkt:implement-primitives} is one-to-one, and it is worth stating in words because the
mapping is cleaner than the seven nominal rates suggest. The emission tax $\tau^\Xi=p^P$ prices the
pollution stock. The primary-waste charge $\tau^{\circ}=p^W$ prices the ledger: a tonne released from
the capital stock, or displaced from nature, adds a tonne to the flow the economy must handle, and it
pays the full social cost of that tonne. The use-waste charge $\tau^{\mathcal D}=\zeta$ prices the
intensity condition: a tonne diffused when a final good is used up does \emph{not} add a tonne to the
flow, because the total is pinned by $R$ and by what gross investment stores, so it pays only the
composition price $\zeta=\sigma\left(p^W-p^M\right)$. Everything else in the instrument list is a
consequence of these three.

\smalltitle{Why the two gate fees differ, and by how much.} The gap $\tau^{\circ}-\tau^{\mathcal D}=p^W-\zeta$
is precisely the materials input tax $\tau^R$, and this is not a coincidence but an accounting identity.
A tonne entering production ends up in one of two places: diffused as use waste, where the gate fee
collects only $\zeta$, or stored in new capital, where it is rebated at $\tau^R$ and later pays
$\tau^{\circ}=p^W$ on release. Charging $\tau^R$ at the point of entry and rebating it on the stored
fraction tops the first route up from $\zeta$ to $p^W$ and leaves the second untouched. That is why the
three rates $\tau^{\mathcal D}$, $\tau^{\circ}$ and $\tau^R$ satisfy one linear restriction and why
Corollary~\ref{cor:mkt:budget} comes out exactly balanced. A tonne of material pays $p^W$ in total,
once, however it travels; the instruments only decide where along the route it is collected.

\smalltitle{A materials deposit--refund.} Block~(B) has a familiar name in the policy literature. Every
tonne drawn into production carries a deposit $\tau^R$, and the deposit is refunded on the material
that is locked into durable capital rather than diffused. The scheme is what makes the market's
investment margin agree with the planner's: without the refund the household would face
$q=1+\Phi^Ip^M$ and would treat embodied material as a pure liability, whereas the planner's
$q=1-\Phi^I\left(1-\sigma\right)\left(p^W-p^M\right)$ is typically \emph{below} unity, because
investing holds mass out of today's waste flow. The refund converts the private liability into the
social premium. The sign of the intervention on this margin follows from the identity
$\tau^R=p^W-\zeta=\left(1-\sigma\right)p^W+\sigma p^M$: the deposit is a convex combination of the two
material shadow costs, so it is positive whenever both are, and durable capital formation is then
subsidised at the margin relative to laissez-faire. Proposition~\ref{prop:sp:phases}(ii) makes this
unconditional in the semi-linear phase, where $p^W>0$ throughout and hence $p^M>0$
by~\eqref{eq:sp:sum:M}. It can fail only in the circular phase, and only once waste has turned into a
resource; whether the deposit ever turns into a levy on investment is one of the sign questions left
open in Section~\ref{sec:sp:opt:results}.

\smalltitle{What the government has to know.} The rates in~\eqref{eq:mkt:implement-derived} are
functions of time and of the first-best path, as Pigouvian rates always are. Two of them additionally
require the government to know the material intensities: $s^I$ needs $\Phi^I$, and the gate fees
operate on tonnages $\Omega^jx_j$ that have to be measured activity by activity. This is a real
informational requirement, but it is the one thing in the model that existing statistics are actually
built to deliver. As noted in Section~\ref{sec:sp:setup:aggregatematerial}, $\Omega$ and the relative
intensities $\omega^j$ exist in this model precisely as the bridge between physical flows and reported
material-flow accounts; here that bridge is load-bearing, since the tax base is a tonnage and the
accounts are what measure it.

\smalltitle{Ad valorem equivalents.} It may be more natural, for comparison with the tax systems
actually in use, to state the gate fee as a set of taxes on activities rather than on tonnages. Since
the waste generated per unit of activity $j$ is $\Omega^j=\Omega\omega^j$, the charge
$\tau^{\mathcal D}=\zeta$ is equivalent to the schedule
\begin{align}
\tau^C &= \zeta\Omega^C,
&
\tau^Y &= \zeta\Omega^Y,
&
\tau^N &= \zeta\Omega^N,
&
\tau^D &= \zeta\Omega^D,
&
\tau^{c^W} &= \zeta\Omega^W,
\label{eq:mkt:ad-valorem}
\end{align}
namely a consumption tax, an output tax, and ad valorem taxes on extraction effort, exploration effort
and treatment expenditure, each at a rate equal to the shadow value of embodied material times the
material intensity of the activity. Five instruments collapse into one because they share a base. The
reduction is the model's most direct policy statement: what should be taxed is not the activity but
the material it diffuses, and once the tax is levied on that base the differentiation across
activities takes care of itself.

\smalltitle{What is \emph{not} required.} Two omissions deserve comment. First, exploration bears no
tax beyond the gate fee on the mass it displaces, even though it is an environmentally costly activity
in the ordinary sense. The reason is the one recorded in Section~\ref{sec:sp:opt:foc}: discovery adds
nothing to $R$, so it draws nothing through the ledger, and it is a complement to extraction rather
than a substitute for it in the materials accounting. Second, recycling receives no subsidy for
recycling as such. It receives $s^R=d^Wp^P$, which is a pollution subsidy --- the leakage the recovered
tonne avoids --- and nothing more. The materials value of recycling is delivered entirely through the
price $P^R$ that the deposit $\tau^R$ has already raised at the mine. This is the market counterpart of
\emph{recycling is paid the cost of the virgin tonne it displaces, and nothing else}: in the planned
economy the statement is about $V$; here it is about the price a recycler actually receives.

\subsection{Phases in the market economy}\label{sec:mkt:implementation:phases}

Because Proposition~\ref{prop:mkt:implementation} matches the two Kuhn--Tucker systems branch by
branch, the phase structure of Proposition~\ref{prop:sp:phases} carries over intact, and it acquires an
observable marker that the planned economy does not have.

\begin{corollary}[The market for recyclable waste dates the transition]\label{cor:mkt:phases}
Under the policy of Proposition~\ref{prop:mkt:implementation}, the regulated market economy passes
through the same two phases as the planned economy, at the same date $t_R$. In the semi-linear phase
the price of treated waste is zero, $P^T=0$, and the waste sector treats waste solely because
treatment converts a tonne from the rate $\tau^U=p^P$ to the rate $\tau^T=d^Wp^P$, so that
$\varpi$ solves $p^P\left(1-d^W\right)=c^{T\prime}\!\left(\varpi\right)\left(1+\zeta\Omega^W\right)$.
The price $P^T$ becomes positive exactly at $t_R$ and rises thereafter.
\end{corollary}

\begin{proof}
By Proposition~\ref{prop:mkt:implementation} the equilibrium allocation is the first best, so the
phases and $t_R$ are those of Proposition~\ref{prop:sp:phases}. In the semi-linear phase $K^R=0$, so
$\alpha=0$ and $P^T=\alpha V=0$ by~\eqref{eq:mkt:recycler-T}; equivalently, recyclers demand no treated
waste at any positive price, and free disposal clears the market at zero. The treatment
condition~\eqref{eq:mkt:waste-foc} then reduces to the stated equality, which is
Proposition~\ref{prop:sp:phases}(iii) written with market instruments. Continuity of $\alpha$ at $t_R$
--- part~(iv) of that proposition --- gives continuity of $P^T$ at zero.
\end{proof}

The corollary is worth more than its proof. In the planned economy the onset of the circular phase is
a statement about shadow prices and is not directly observable. In the market economy it is the date
at which waste acquires a positive market value, which is a fact about an actual price series. It also
sharpens what the semi-linear phase means institutionally: a waste sector that sorts and treats without
any recycling taking place is not a paradox, it is the efficient response to a regulation that
distinguishes controlled disposal from uncontrolled disposal. That is a recognisable description of
waste policy in most countries for most of the twentieth century.

\subsection{The unregulated economy}\label{sec:mkt:implementation:unregulated}

Setting every instrument to zero gives the polar case, and it is worth recording what survives, because
the paper's quantitative question is how far and for how long the unregulated economy departs from the
optimum. With $\tau^\Xi=\tau^{\mathcal D}=\tau^{\circ}=0$, hence $\tau^R=s^I=s^R=0$, the private
conditions collapse to
\begin{align}
F_R &= P^R,
&
F_K &= P^K,
&
P^R &= C^N_N+P^S,
&
a'\!\left(x\right)P^R &= P^K,
&
\alpha P^R &= P^T,
&
c^{T\prime}\!\left(\varpi\right) &= P^T,
\label{eq:mkt:unregulated}
\end{align}
with $q=1$ and $s^W=c^W-P^T\varpi$, which the tender can only deliver as a positive subsidy if the
government funds it; absent a collection obligation altogether, waste is collected only where
$P^T\varpi\ge c^W$, and is otherwise released.

Three observations follow immediately, and a fourth does not.

First, the materials side of the model disappears entirely. Every wedge $1\pm\zeta\Omega^j$ is gone,
the ledger is unpriced, and $M^K$ ceases to be an object anyone tracks. The unregulated economy
behaves, on these margins, exactly like a model with no materials accounting at all --- which is a
useful check on the framework rather than a defect of it.

Second, extraction is unambiguously too cheap. The optimal condition charges the miner
$C^N_N\zeta\Omega^N+p^W\Omega^{N,S}$ per tonne over and above the private marginal cost, both terms
strictly positive whenever $\zeta>0$ and $p^W>0$. At any given $P^R$ and $P^S$ the unregulated firm
extracts more.

Third, and this is the substantive one, recycling starts \emph{later}. Compare the corner
branch~\eqref{eq:mkt:recycler-K-corner} in the two economies: the regulated recycler holds
$a'\!\left(0^+\right)\left(P^R+d^Wp^P\right)$ against $P^K$, the unregulated one holds
$a'\!\left(0^+\right)P^R$ against $P^K$. At common prices the unregulated threshold is harder to
clear, by exactly the leakage subsidy, so the semi-linear phase lasts longer. The comparison is a
partial one --- $P^R$, $P^K$ and the whole path of the reserve differ across the two economies, so it
bounds nothing on its own --- but it identifies the mechanism, and it identifies which term is
responsible.

What does \emph{not} follow is the magnitude, or whether the unregulated economy ever reaches the
circular phase at all. Both sides of the threshold move with development: depletion raises $C^N_N$ and
$P^S$ and therefore $P^R$, while capital accumulation lowers $P^K$ and productivity growth raises it.
Section~\ref{sec:sp:opt:results} records the same non-monotonicity for the planner. Whether the
unregulated transition merely happens late or fails to happen, and what the delay costs in welfare
terms, is a quantitative question, and it is one of the questions the numerical model in the next part
is built to answer.
